\section{Introduction}
\label{sec:intro}

Every variational quantum algorithm has the same shape as a neural network: a fixed computational graph (the circuit) with continuous parameters (the gate angles), a loss computed by running the graph, and a classical optimizer that adjusts the parameters \citep{cerezo2021variational}. The optimizer is the expensive part. Each loss evaluation is a run of the quantum device, gradients vanish exponentially in the number of qubits for generic circuits \citep{mcclean2018barren}, and the optimization problem is NP-hard in general \citep{bittel2021training}. A great deal of work therefore goes into choosing good initial parameters, or into reusing parameters across problem instances.

This paper asks a simple question: \emph{can a large language model write the parameters directly?} Not propose a circuit for someone else to train, and not act as an optimizer that is shown scores and asked to improve them, but read a description of the problem---a Hamiltonian, a graph, ten pictures of handwritten digits---and reply with the numbers that go into the gates, with no feedback of any kind. \citet{takehi2026sparks} recently showed that a frontier model can do this for classical neural networks: read the weights of a trained network and say what it does, edit them to fix a failing policy, and write the weights of game-playing networks from scratch. They call it \emph{weight-space programming}---a general-purpose LLM reads and writes a network's weights directly, not through an optimizer, around a natural-language specification and sometimes with no environment feedback at all. Our work is inspired by theirs and carries the idea to a domain where it has a particular pull: variational quantum circuits have few parameters, each with a physical or geometric meaning, and the cost of the optimizer they replace is paid in runs of a quantum device.

We test the question on the three main families of variational quantum algorithms: supervised classification with a parameterized circuit (QML), ground-state energy estimation (VQE), and combinatorial optimization (QAOA). In every experiment the model receives the problem, the exact gate conventions and the circuit (or, in one QML setting, a small circuit language in which it may design its own), and returns a JSON list of angles. The angles are simulated exactly as written. We compare them against three references on the same circuit: random angles, a classical optimizer started from random angles, and the same optimizer started from the model's angles. All experiments use one frontier model, \gpt, and every prompt, reply and score is released.

\paragraph{Findings.}
\begin{enumerate}
  \item \textbf{Weight programming works when the weights have meaning.} With the Hamiltonian variational ansatz for the transverse-field Ising chain, the model's angles capture 95--100\% of the spectrum between the highest and lowest eigenvalue on every problem, equal to the optimizer; for H$_2$ it derives the energy as a function of the single angle and writes its exact minimizer; for QAOA at depth one it derives the closed-form expectation of the cut and maximizes it analytically, matching the optimizer on all six graphs. For digit classification it hand-writes 12--18 weights of a fixed amplitude-encoding circuit and scores 93--96\% on 3 vs 8, \emph{above} an optimizer trained on the same ten examples.
  \item \textbf{The model's weights are a strong warm start.} Started from the model's angles, the optimizer reaches its best value after a median of one loss evaluation in VQE and QAOA, against hundreds from random starts, and it reaches that value more often.
  \item \textbf{Left to design the circuit, the model picks the one it can reason about.} In the free-design QML setting the model converges, in 30 of 30 trials across six conditions, on chains of data-dependent rotations whose weights are pixel importances: a single-qubit chain for binary tasks (a linear classifier in disguise) and one such chain per qubit, read out compositionally, for ten classes, with hand-chosen weights that score 91--95\% on binary tasks and 65\% on all ten digits. When structural rules forbid the single-qubit chain, it finds circuit identities that compute the same weighted sum on two qubits, or writes a nearest-prototype classifier by hand. Interpretability of the weight space, not quantumness, is what the model exploits.
  \item \textbf{It does not scale to opaque weight spaces.} When the fixed circuit is deepened, so that each weight's effect on the output can no longer be reasoned about, accuracy on 3 vs 8 falls from 0.96 at 18 weights to 0.70 at 54, and on ten classes it never leaves chance (0.07--0.14) while the optimizer on the same examples climbs to 0.57. The failure is informative: the model says it is writing ``heuristic weights'' and it is right.
  \item \textbf{Conventions are a failure mode.} In 6 of 42 QAOA trials the model produced the correct angle magnitudes with every mixer angle negated; those circuits score near zero, and with the sign corrected they match the optimizer. Weight programming is only as good as the shared understanding of what a weight means.
\end{enumerate}

These results extend weight-space programming from classical networks to quantum circuits and reframe what a language model can contribute to variational quantum algorithms: not as a replacement for the optimizer in general, but as a source of parameters for circuits whose weights carry physical or geometric meaning, and as an initializer that removes most of the optimizer's work.
